\begin{examplebox}
\textbf{例 1：估算 $\sin0.2$，四舍五入到小数点后四位。}
这里 $0.2$ 是弧度。仅说“$x$ 小”不能保证最后一位；应先算多项式，再给出误差。
\[
 S_5(0.2)=\frac15-\frac{(1/5)^3}{6}+\frac{(1/5)^5}{120}
        =\frac{74501}{375000}=0.1986693333\ldots
\]
由严格误差界，
\[
 |\sin0.2-S_5(0.2)|\le\frac{0.2^7}{5040}<2.54\times10^{-9}.
\]
所以真值落在 $(0.1986693307,\,0.1986693359)$ 内。
该区间完全位于会舍入成 $0.1987$ 的区间
$[0.19865,0.19875)$ 中，故答案为 $\boxed{0.1987}$。
用更高精度独立计算，$\sin0.2=0.198669330795\ldots$，与误差界一致。

下图记三次近似为 $S_3(x)=x-x^3/6$，五次近似仍为 $S_5(x)$。
\begin{center}
\begin{tikzpicture}[font=\small]
\draw[black!30,rounded corners=2pt] (0,0) rectangle (14.4,8.2);
\node[anchor=west,font=\small\bfseries] at (0.5,7.72) {正弦估值：同一横坐标下比较三次与五次近似误差};
\begin{axis}[width=11.2cm,height=6.1cm,at={(2.05cm,1.25cm)},anchor=south west,
axis lines=left,xlabel={$x$},ylabel={$y$},tick label style={font=\small},label style={font=\small},
legend style={font=\small,draw=none,fill=white,at={(0.03,0.97)},anchor=north west},
legend cell align=left,unbounded coords=jump,xmin=0.1,xmax=1,ymin=1e-12,ymax=1e-2,ymode=log,ylabel={绝对误差},restrict x to domain=0.1:1,legend columns=3,legend style={draw=none,fill=none,at={(0.5,1.14)},anchor=south,column sep=3pt}]
\addplot[purple,thick,dashdotted] table[x=x,y=es3] {assets/tikz/F004_samples.dat};
\addlegendentry{$|\sin x-S_3(x)|$}
\addplot[orange!85!black,thick,dashed] table[x=x,y=es5] {assets/tikz/F004_samples.dat};
\addlegendentry{$|\sin x-S_5(x)|$}
\addplot[black,densely dotted,domain=0.1:1] {0.00005};\addlegendentry{$5\times10^{-5}$}
\end{axis}

\end{tikzpicture}
\end{center}

图中误差来自真实函数与各多项式之差，纵轴为对数尺度。
虚点线是半个 $0.0001$；真正保证四舍五入仍须像本题一样，检查整个误差区间没有跨过舍入分界。

\textbf{例 2：抵消后的首个非零项决定极限。}
求 $\displaystyle\lim_{x\to0}\frac{\cos x-1+x^2/2}{x^4}$。
\[
\begin{aligned}
 \cos x-1+\frac{x^2}{2}&=\frac{x^4}{24}+R(x),\\
 |R(x)|&\le\frac{|x|^6}{720}.
\end{aligned}
\]
对 $x\ne0$，除以 $x^4$ 后的误差满足
\[
 \left|\frac{\cos x-1+x^2/2}{x^4}-\frac1{24}\right|
 \le\frac{x^2}{720}\longrightarrow0.
\]
故极限为 $\boxed{1/24}$。这一不等式同时回验左右极限，无须把一个很小的浮点差值当成精确的零。

\begin{center}
\begin{tikzpicture}[font=\small]
\draw[black!30,rounded corners=2pt] (0,0) rectangle (14.4,8.2);
\node[anchor=west,font=\small\bfseries] at (0.5,7.72) {抵消后再除以 $x^4$：留下的系数是 $1/24$};
\begin{axis}[width=11.2cm,height=6.1cm,at={(2.05cm,1.25cm)},anchor=south west,
axis lines=left,xlabel={$x$},ylabel={$y$},tick label style={font=\small},label style={font=\small},
legend style={font=\small,draw=none,fill=white,at={(0.03,0.97)},anchor=north west},
legend cell align=left,unbounded coords=jump,xmin=-0.8,xmax=0.8,ymin=0.0406,ymax=0.0421,scaled y ticks=false,yticklabel style={/pgf/number format/fixed,/pgf/number format/precision=4},legend columns=2,legend style={draw=none,fill=none,at={(0.5,1.14)},anchor=south,column sep=8pt}]
\addplot[blue,very thick] table[x=x,y=cr] {assets/tikz/F004_samples.dat};
\addlegendentry{$(\cos x-1+x^2/2)/x^4$}
\addplot[purple,densely dotted,domain=-0.8:0.8] {1/24};\addlegendentry{$1/24$}
\addplot[only marks,mark=o,mark options={fill=white},blue,forget plot] coordinates {(0,0.0416666666667)};
\end{axis}

\end{tikzpicture}
\end{center}

空心点表示原商在 $x=0$ 无定义，而左右都趋向 $1/24$。
若只用 $\cos x\approx1-x^2/2$，图中的整条残差曲线就会被错误地抹去。

\textbf{例 3：复合函数的二次近似与另一方法回验。}
求 $f(x)=\ln(1+\sin x)$ 在零点附近的二次近似。
取 $|x|\le1/2$，由 $|\sin x|\le|x|$ 得 $1+\sin x\ge1/2>0$，定义域安全。
设 $u=\sin x$，则 $u=x+O(x^3)=O(x)$，且 $u^2=x^2+O(x^4)$。
\[
\begin{aligned}
 \ln(1+u)&=u-\frac{u^2}{2}+O(u^3),\\
 \ln(1+\sin x)&=x-\frac{x^2}{2}+O(x^3).
\end{aligned}
\]
答案为 $\boxed{x-x^2/2}$。
\textbf{回验：} 直接求导得
\[
 f(0)=0,\quad f'(x)=\frac{\cos x}{1+\sin x},\quad
 f''(x)=-\frac1{1+\sin x}.
\]
因而 $f'(0)=1,f''(0)=-1$，且 $f$ 在零点附近三阶连续可导，确实得到相同系数与 $O(x^3)$ 余项。

\begin{center}
\begin{tikzpicture}[font=\small]
\draw[black!30,rounded corners=2pt] (0,0) rectangle (14.4,8.2);
\node[anchor=west,font=\small\bfseries] at (0.5,7.72) {复合函数：先代入，再保留到二次项};
\begin{axis}[width=11.2cm,height=6.1cm,at={(2.05cm,1.25cm)},anchor=south west,
axis lines=left,xlabel={$x$},ylabel={$y$},tick label style={font=\small},label style={font=\small},
legend style={font=\small,draw=none,fill=white,at={(0.03,0.97)},anchor=north west},
legend cell align=left,unbounded coords=jump,xmin=-0.6,xmax=0.6,ymin=-1,ymax=0.6,legend style={at={(0.03,0.97)},anchor=north west}]
\addplot[blue,very thick] table[x=x,y=composite] {assets/tikz/F004_samples.dat};
\addlegendentry{$\ln(1+\sin x)$}
\addplot[orange!85!black,thick,dashed] table[x=x,y=q2] {assets/tikz/F004_samples.dat};
\addlegendentry{$x-x^2/2$}

\end{axis}

\end{tikzpicture}
\end{center}

\end{examplebox}
